\section{Introdu\c{c}\~ao}
\label{sec:introducao}

% O QUE INCLUIR:
% - Contexto e motivacao (IA aplicada a esfera publica / debates legislativos).
% - O problema especifico que seu grupo escolheu atacar e por que ele importa.
% - A trilha de inspiracao adotada (A, B, C, D) ou a tarefa inedita proposta.
% - Uma sintese das suas contribuicoes, em lista, ao final desta secao.

<Apresente o contexto do desafio e o problema escolhido pelo grupo.>
A solu\c{c}\~ao proposta neste trabalho pode ser citada e referenciada ao longo
do texto; refer\^encias usam o comando \verb|\citep{}| ou \verb|\citet{}|, por
exemplo \citep{lewis2020rag}.

\medskip\noindent\textbf{Contribui\c{c}\~oes.} As principais contribui\c{c}\~oes deste trabalho s\~ao:
\begin{itemize}
  \itemsep2pt
  \item <Contribui\c{c}\~ao 1 --- ex.: nova arquitetura / m\'etodo / an\'alise.>
  \item <Contribui\c{c}\~ao 2 --- ex.: protocolo de avalia\c{c}\~ao proposto.>
  \item <Contribui\c{c}\~ao 3 --- ex.: resultados, dataset derivado ou ferramenta.>
\end{itemize}
